\documentclass[11pt]{article}
\usepackage[localtheorems]{raeez-math-template}

%%% Packages

%%% Page geometry

%%% Theorem environments
\theoremstyle{plain}
\newtheorem{theorem}{Theorem}[section]
\newtheorem{proposition}[theorem]{Proposition}
\newtheorem{corollary}[theorem]{Corollary}
\newtheorem{lemma}[theorem]{Lemma}
\newtheorem{conjecture}[theorem]{Conjecture}

\theoremstyle{definition}
\newtheorem{definition}[theorem]{Definition}
\newtheorem{example}[theorem]{Example}

\theoremstyle{remark}
\newtheorem{remark}[theorem]{Remark}

%%% Macros
\newcommand{\fg}{\mathfrak{g}}
\newcommand{\fh}{\mathfrak{h}}
\newcommand{\fsl}{\mathfrak{sl}}
\newcommand{\cA}{\mathcal{A}}
\newcommand{\cW}{\mathcal{W}}
\newcommand{\Mbar}{\overline{\mathcal{M}}}
\newcommand{\Res}{\operatorname{Res}}
\newcommand{\Vir}{\mathrm{Vir}}
\newcommand{\barBch}{\bar{B}^{\mathrm{ch}}}
\DeclareMathOperator{\ad}{ad}
\DeclareMathOperator{\End}{End}
\DeclareMathOperator{\Conf}{Conf}
\DeclareMathOperator{\Bun}{Bun}
\DeclareMathOperator{\Tw}{Tw}
\numberwithin{equation}{section}

%%% Title
\title{The genus-$1$ seven faces:\\
elliptic collision residues and the KZB connection}
\author{Raeez Lorgat}
\date{}

\begin{document}
\maketitle

%%% ================================================================
%%% ABSTRACT
%%% ================================================================

\begin{abstract}
At genus~$0$, the collision residue of a chirally Koszul chiral
algebra admits seven equivalent presentations: as a twisting
morphism, a line-operator $r$-matrix, a $\lambda$-bracket, commuting
sphere Hamiltonians, the Drinfeld $r$-matrix, the Sklyanin bracket,
and the Gaudin Hamiltonians. We prove the genus-$1$ extension: on an
elliptic curve $E_\tau$, the genus-$1$ collision residue
$r_\cA^{(1)}(z,\tau)$ is the elliptic regularization of the genus-$0$
residue, obtained by replacing rational propagator kernels with
prime-form elliptic kernels. All seven faces at genus~$1$ specialize a
\emph{single $E_1$-datum} $\Theta_\cA$ via the dictionary
$\Psi \rightsquigarrow \hbar \rightsquigarrow \tau$: $\Psi$ is the
algebraic spectral parameter on genus-$0$ faces, $\hbar$ is the
quantisation parameter on the Drinfeld / Sklyanin / Yangian faces,
and $\tau$ is the elliptic modulus on the KZB / Belavin--Drinfeld /
Felder faces. The passage is the bar-intrinsic projection
$\Theta_\cA^{(g=1)} = \mathrm{Reg}_\tau \circ \Theta_\cA^{(g=0)}$
where $\mathrm{Reg}_\tau$ is elliptic regularisation.
For affine Kac--Moody algebras $\cA = \widehat{\fg}_k$, the seven
faces specialize to the Knizhnik--Zamolodchikov--Bernard connection,
the Belavin--Drinfeld elliptic $r$-matrix, the Felder dynamical
elliptic $R(z,\tau)$ (via Fay trisecant), and the elliptic Gaudin
Hamiltonians. The prime-form $B$-cycle shift is the local source of
the Drinfeld--Kohno quantum-group parameter; with
\(\hbar=(k+h^\vee)^{-1}\), the standard parameter is
\(q_{\mathrm{DK}}=\exp(\pi i\hbar)\) and the full-cycle factor is
\(q_{\mathrm{DK}}^2\). For class~$M$ algebras (Virasoro,
$\cW_N$), the collision residue produces higher-order elliptic
differential operators involving \(P_\tau,P_\tau',\ldots,P_\tau^{(2N-3)}\)
outside the affine Hitchin--KZB $r$-matrix lane.
\end{abstract}

\begin{remark}[Genus-$1$ seven-face datum]\label{rem:genus1-seven-face-datum}
The genus-$1$ comparison kernels are obtained by restricting the
bar-intrinsic Maurer--Cartan element to the elliptic collision
stratum. The KZB, Belavin--Drinfeld, Felder, and Gaudin readouts are
different projections of this restricted datum. The prime-form
normalization uses
\[
\xi_\tau(z)=\partial_z\log\theta_1(z\mid\tau)
=\zeta_\tau(z)-2\eta_1z,\qquad
P_\tau(z)=-\partial_z\xi_\tau(z)=\wp(z,\tau)+2\eta_1.
\]
The quasi-period \(\xi_\tau(z+\tau)-\xi_\tau(z)=-2\pi i\) is the
local monodromy input; the quantum-group parameter also requires the
chosen Drinfeld--Kohno or Kazhdan--Lusztig normalization.
\end{remark}


%%% ================================================================
%%% 1. INTRODUCTION
%%% ================================================================

\section{Introduction}\label{sec:intro}

The genus-$0$ seven-face theorem~\cite{Lorgat26} identifies the
collision residue $r_\cA(z) = \Res^{\mathrm{coll}}_{0,2}(\Theta_\cA)$
of a chirally Koszul chiral algebra $\cA$ with seven classical
objects: a bar-cobar twisting morphism, a line-operator $r$-matrix in
the Dimofte--Niu--Py framework~\cite{DNP25}, a classical
$\lambda$-bracket in the Khan--Zeng 3d sigma-model~\cite{KhanZeng25},
the Gaiotto--Zeng commuting sphere Hamiltonians, the Drinfeld
$r$-matrix, the Sklyanin Poisson bracket, and the Gaudin
Hamiltonians. The seven presentations agree as elements of
$\End_\cA(2)[\![z^{-1}]\!]$; their mutual identification is
controlled by the collision depth $k_{\max}$ and the shadow class
$G/L/C/M$.

The collision residue $r_\cA(z)$ is the degree-$2$ component of
the $E_1$ Maurer--Cartan element
$\Theta^{E_1}_\cA$ in the ordered convolution algebra
$\mathfrak{g}^{E_1}_\cA$, where
$B^{\mathrm{ord}}(\cA) = T^c(s^{-1}\bar{\cA})$ is the cofree tensor
coalgebra with deconcatenation coproduct. Its
$\Sigma_n$-coinvariant projection is the scalar shadow $\kappa$:
$\mathrm{av}(r(z))$ in abelian and scalar families, and
$\mathrm{av}(r(z))+\dim(\fg)/2$ for non-abelian affine
Kac--Moody;
the spectral parameter that $\kappa$ discards is precisely the data
the seven faces read. At genus~$1$, the spectral parameter acquires
modular dependence: $r^{(1)}_\cA(z,\tau)$ depends on the elliptic
modulus $\tau$, and the $B$-cycle monodromy of the ordered collision
residue encodes the quantum group parameter.

This paper extends the identification to genus~$1$. The elliptic
curve $E_\tau = \mathbb{C}/(\mathbb{Z}+\mathbb{Z}\tau)$ replaces the
Riemann sphere, the rational propagator $1/z$ is replaced by the
prime-form kernel
\(\xi_\tau(z)=\theta_1'(z,\tau)/\theta_1(z,\tau)\), and each rational
function $1/z^{n+1}$ in the collision expansion is replaced by the
corresponding derivative of
\(P_\tau(z)=-\partial_z\xi_\tau(z)\). The resulting genus-$1$
collision residue
$r_\cA^{(1)}(z,\tau)$ admits seven equivalent presentations that
degenerate to the genus-$0$ faces as $\tau \to i\infty$.

The main results are:
\begin{enumerate}[label=(\arabic*)]
\item The \emph{elliptic regularization theorem}
(Theorem~\ref{thm:elliptic-reg}): the genus-$1$ collision residue is
\[
r_\cA^{(1)}(z,\tau)
\;=\;
c_0\,\xi_\tau(z)
\;+\;
\sum_{n \geq 1}
\frac{(-1)^{n-1}}{n!}\,c_n \cdot P_\tau^{(n-1)}(z),
\]
where $c_n = a_{(n)}b$ are the OPE-mode collision coefficients.

\item The identification of the $dz$-component of the \emph{KZB
connection} with the genus-$1$ collision residue, and of the
$d\tau$-component with its $z$-derivative
(Theorem~\ref{thm:kzb}).

\item The identification of the genus-$1$ collision residue with
the \emph{Belavin--Drinfeld elliptic $r$-matrix}
(Theorem~\ref{thm:elliptic-r}).

\item The \emph{elliptic Gaudin comparison}: the two-point and Cartan
sectors match the KZB kernel; full \(n\geq3\) flatness is the Felder
dynamical Yang--Baxter statement (Theorem~\ref{thm:gaudin}).

\item The \emph{class~$M$ collision residue} for Virasoro and
$\cW_N$: higher-order elliptic operators outside the affine
Hitchin--KZB \(r\)-matrix lane
(Theorems~\ref{thm:virasoro-g1} and~\ref{thm:wn-g1}).

\item The \emph{degeneration theorem}: all seven genus-$1$ faces
degenerate continuously to their genus-$0$ counterparts as
$q = e^{2\pi i\tau} \to 0$ (Theorem~\ref{thm:degeneration}).
\end{enumerate}

The prime-form kernel defines the elliptic collision residue. Its
\(z\)- and \(\tau\)-derivatives give the KZB connection, its
root-channel continuation gives the Belavin--Drinfeld kernel, and its
collision-depth projections give the Gaudin and class-\(M\) operators.


%%% ================================================================
%%% 2. THE ELLIPTIC COLLISION RESIDUE
%%% ================================================================

\section{The elliptic collision residue}\label{sec:propagator}

\subsection{The genus-$1$ bar propagator}

At genus~$0$ the bar propagator is $\eta_{12} = d\log(z_1 - z_2)$: a
meromorphic one-form on $\mathbb{P}^1$ with a simple pole along the
diagonal. On the elliptic curve $E_\tau$, Liouville's theorem forbids
a meromorphic function with a single simple pole that is doubly
periodic. The propagator must be replaced by a quasi-periodic object.

\begin{definition}[Genus-$1$ bar propagator]\label{def:bar-propagator}
The \emph{genus-$1$ bar propagator} on $E_\tau$ is the logarithmic
derivative of the prime form:
\begin{equation}\label{eq:propagator}
\eta_{12}^{E_\tau}
\;=\;
d\log E(z_1,z_2)
\;=\;
\xi_\tau(z_1 - z_2)\,d(z_1 - z_2),
\end{equation}
where $E(z_1,z_2) = \theta_1(z_1 - z_2 \mid \tau)/\theta_1'(0\mid\tau)$
is the prime form on $E_\tau$ (a section of
$K^{-1/2} \boxtimes K^{-1/2}$), $\theta_1$ is the Jacobi theta
function, and
\[
\xi_\tau(z):=\partial_z\log\theta_1(z\mid\tau)
=\frac{\theta_1'(z\mid\tau)}{\theta_1(z\mid\tau)}
=\zeta_\tau(z)-2\eta_1z
\]
is the prime-form logarithmic kernel. Here \(\zeta_\tau\) is the
Weierstrass zeta function for periods \((1,\tau)\):
\begin{equation}\label{eq:weierstrass-zeta}
\zeta_\tau(z)
\;=\;
\frac{1}{z}
+ \sum_{(m,n) \neq (0,0)}
\biggl[
 \frac{1}{z - m - n\tau}
 + \frac{1}{m + n\tau}
 + \frac{z}{(m + n\tau)^2}
\biggr].
\end{equation}
\end{definition}

Put \(P_\tau(z):=-\partial_z\xi_\tau(z)=\wp(z,\tau)+2\eta_1\).
The prime-form kernel has Laurent expansion
$\xi_\tau(z) = 1/z + O(z)$ near $z = 0$, matching the genus-$0$
propagator $1/z$ at leading order. The difference is global:
$\xi_\tau$ is quasi-periodic,
\begin{equation}\label{eq:quasi-periodicity}
\xi_\tau(z+1) = \xi_\tau(z),
\qquad
\xi_\tau(z+\tau) = \xi_\tau(z) - 2\pi i.
\end{equation}
The quasi-periodicity of the propagator is the geometric origin of the
$B$-cycle monodromy that distinguishes genus~$1$ from genus~$0$ in
all seven faces.

\begin{remark}[Propagator weight]\label{rem:propagator-weight}
The bar propagator $d\log E(z,w)$ has weight~$1$ in both variables,
regardless of the conformal weight of the fields being sewed. The
logarithmic derivative of the prime form is a $(1,0)$-form in each
variable; the conformal weight $h$ of the generator enters the
OPE-mode coefficient $c_n$, not the propagator kernel.
\end{remark}


\subsection{The collision residue}

\begin{definition}[Genus-$1$ collision residue]\label{def:collision-residue}
For a chirally Koszul chiral algebra $\cA$ on $E_\tau$, the
\emph{genus-$1$ collision residue} is
\begin{equation}\label{eq:collision-residue}
r_\cA^{(1)}(z,\tau)
\;:=\;
\Res^{\mathrm{coll}}_{1,2}(\Theta_\cA)\big|_{E_\tau}
\;\in\;
\End_\cA(2)[\![z^{-1}]\!] \otimes \mathbb{C}[\![\tau]\!],
\end{equation}
where $\Theta_\cA$ is the bar-intrinsic Maurer--Cartan element.
\end{definition}

\begin{theorem}[Elliptic regularization]\label{thm:elliptic-reg}
Let $\cA$ be a chirally Koszul chiral algebra with genus-$0$ collision
residue $r_\cA(z) = \sum_{n \geq 0} c_n/z^{n+1}$, where
$c_n = a_{(n)}b$ are the OPE-mode coefficients. The genus-$1$
collision residue is
\begin{equation}\label{eq:elliptic-expansion}
r_\cA^{(1)}(z,\tau)
\;=\;
c_0\,\xi_\tau(z)
\;+\;
\sum_{n \geq 1}
\frac{(-1)^{n-1}}{n!}\,c_n \cdot P_\tau^{(n-1)}(z),
\end{equation}
where \(P_\tau^{(m)}=d^mP_\tau/dz^m\). Since
\(P_\tau=\wp+2\eta_1\), one has \(P_\tau^{(m)}=\wp^{(m)}\) for
\(m\geq1\). In particular, the first three
terms are
\begin{equation}\label{eq:first-three}
r_\cA^{(1)}(z,\tau)
\;=\;
c_0\,\xi_\tau(z)
\;+\;
c_1\,P_\tau(z)
\;-\;
\tfrac{1}{2}\,c_2\,P_\tau'(z)
\;+\;\cdots
\end{equation}
The Laurent expansion of $r_\cA^{(1)}$ near $z = 0$ agrees with
$r_\cA(z)$ at leading order.
\end{theorem}

\begin{proof}
Each rational function $1/z^{n+1}$ ($n \geq 0$) in the genus-$0$
collision expansion is the unique meromorphic function on
$\mathbb{P}^1$ with a pole of order $n+1$ at $z = 0$ and vanishing
at $z = \infty$. On $E_\tau$, the prime-form normalization fixes the
corresponding polar kernel. For \(n=0\) it is \(\xi_\tau(z)\); for
\(n\geq1\) it is \((-1)^{n-1}P_\tau^{(n-1)}(z)/n!\). The prefactor is
fixed by the Laurent expansion
$P_\tau^{(m)}(z) = (-1)^m(m+1)!/z^{m+2} + O(1)$ at $m = n-1$:
\[
\frac{(-1)^{n-1}}{n!}\,P_\tau^{(n-1)}(z)
\;=\;
\frac{(-1)^{n-1}}{n!}\cdot\frac{(-1)^{n-1}\,n!}{z^{n+1}} + O(z^{-n+1})
\;=\;
\frac{1}{z^{n+1}} + O(z^{-n+1}).
\]
Substituting into the propagator-extraction formula of
Definition~\ref{def:bar-propagator}
yields~\eqref{eq:elliptic-expansion}.
\end{proof}

\begin{remark}[Quasi-periodicity dichotomy]\label{rem:dichotomy}
The kernel \(P_\tau\) and all its derivatives are doubly periodic. The
quasi-periodicity of $r_\cA^{(1)}(z,\tau)$ is therefore inherited
entirely from the \(\xi_\tau\)-term:
\begin{equation}\label{eq:B-cycle-monodromy}
r_\cA^{(1)}(z + \tau,\tau) - r_\cA^{(1)}(z,\tau)
\;=\;
-2\pi i \cdot c_0.
\end{equation}
When $c_0 = 0$ (the Cartan-type case: $\beta\gamma$, lattice VOAs),
the collision residue is a genuine meromorphic function on $E_\tau$.
When $c_0 \neq 0$ (the Yangian-type case: affine Kac--Moody,
Virasoro, $\cW_N$), it is quasi-periodic with non-trivial $B$-cycle
monodromy.
\end{remark}


%%% ================================================================
%%% 3. FACE 4: THE KZB CONNECTION
%%% ================================================================

\section{Face 4: the KZB connection}\label{sec:kzb}

At genus~$0$, Face~4 identifies the collision residue with the
Gaiotto--Zeng commuting sphere Hamiltonians. At genus~$1$, the
sphere Hamiltonians are replaced by the
Knizhnik--Zamolodchikov--Bernard connection on the configuration space
$\Conf_n(E_\tau)$. The KZB connection has two components: a
$dz$-component generalizing the genus-$0$ KZ connection, and a
$d\tau$-component encoding the dependence of conformal blocks on the
modulus $\tau$.

The formulas in this section use the KZ-normalized conformal-block
kernel. In trace-form normalization the binary affine collision
residue is \(k\Omega_{\mathrm{tr}}\xi_\tau(z)\). The KZ coefficient is
\(\Omega_{\mathrm{KZ}}\xi_\tau(z)/(k+h^\vee)\). These are related by
the Sugawara normalization; they are not identical functions of the
level.

\begin{theorem}[Face~4: KZB connection]\label{thm:kzb}
For $\cA = \widehat{\fg}_k$ at level $k \neq -h^\vee$, the modular
shadow connection on $\Conf_n(E_\tau)$ restricts to the KZB
connection:
\begin{equation}\label{eq:kzb-full}
\nabla^{\mathrm{KZB}}
\;=\;
d
\;-\;
\underbrace{
\frac{1}{k+h^\vee}
\sum_{i<j}\Omega_{ij}\,\xi_\tau(z_{ij})\,dz_{ij}
}_{\textup{$dz$-component: genus-$1$ collision residue}}
\;-\;
\underbrace{
\frac{1}{k+h^\vee}
\sum_{i<j}\Omega_{ij}\,P_\tau(z_{ij})\,d\tau
}_{\textup{$d\tau$-component: modular Hamiltonian}},
\end{equation}
where $z_{ij} = z_i - z_j$ and $\Omega_{ij}$ is the Casimir acting
on the $i$-th and $j$-th tensor factors.
\begin{enumerate}[label=\textup{(\roman*)}]
\item The $dz$-component equals the genus-$1$ collision residue:
\begin{equation}\label{eq:kzb-z}
H_{z,i}^{\mathrm{KZB}}
\;=\;
\frac{1}{k+h^\vee}\sum_{j \neq i}\Omega_{ij}\,\xi_\tau(z_{ij})
\;=\;
\sum_{j \neq i}
\bigl[r_{\widehat{\fg}_k}^{(1)}(z_{ij},\tau)\bigr]_{ij}.
\end{equation}
\item The $d\tau$-component equals the $z$-derivative of the collision
residue, through the identity $P_\tau(z)=-\xi_\tau'(z)$:
\begin{equation}\label{eq:kzb-tau}
H_{\tau,i}^{\mathrm{KZB}}
\;=\;
\frac{1}{k+h^\vee}\sum_{j \neq i}\Omega_{ij}\,P_\tau(z_{ij})
\;=\;
-\sum_{j \neq i}
\partial_z\bigl[r_{\widehat{\fg}_k}^{(1)}(z_{ij},\tau)\bigr]_{ij}.
\end{equation}
\end{enumerate}
\end{theorem}

\begin{proof}
\textbf{Part~(i).} The affine Kac--Moody OPE
$J^a(z)\,J^b(w) \sim k\,\delta^{ab}/(z-w)^2 +
f^{ab}_{\;\;c}\,J^c(w)/(z-w)$ has $c_0 = \Omega/(k+h^\vee)$
and $c_n = 0$ for $n \geq 1$: the bar propagator absorbs one pole
order, so the only surviving term is $c_0$. By
Theorem~\ref{thm:elliptic-reg}, the elliptic regularization gives
$r_{\widehat{\fg}_k}^{(1)}(z,\tau) = \Omega\,\xi_\tau(z)/(k+h^\vee)$
in KZ normalization.
Substituting into the Gaiotto--Zeng formula produces~\eqref{eq:kzb-z},
which is the $dz$-component of the KZB
connection~\cite{Bernard88,KZB95}.

\textbf{Part~(ii).} The $d\tau$-component involves the modular
Hamiltonian $D_i$, whose leading term is
$\sum_{j \neq i}P_\tau(z_{ij})\,\Omega_{ij}/(k+h^\vee)$. The
identity \(P_\tau=-\xi_\tau'\) identifies this
with $-\partial_z r_\cA^{(1)}$, yielding~\eqref{eq:kzb-tau}.
\end{proof}

\begin{remark}[Shadow connection interpretation]\label{rem:shadow}
The KZB connection is the genus-$1$ instance of the modular shadow
connection $\nabla^{\mathrm{hol}}_{g,n}$ on $\Mbar_{g,n}$: at
genus~$0$ it gives the KZ connection, at genus~$1$ the KZB
connection, and at genus~$\geq 2$ it produces new flat connections on
configuration spaces of higher-genus curves.
\end{remark}


%%% ================================================================
%%% 4. FACE 5: THE ELLIPTIC r-MATRIX
%%% ================================================================

\section{Face 5: the elliptic $r$-matrix}\label{sec:r-matrix}

At genus~$0$, Face~5 identifies the collision residue with the
KZ-normalized rational kernel $\Omega/((k+h^\vee)z)$. At genus~$1$, the
rational $r$-matrix is replaced by the classical elliptic $r$-matrix
of Belavin~\cite{Belavin81} and
Belavin--Drinfeld~\cite{BelavinDrinfeld82}.

\begin{theorem}[Face~5: elliptic $r$-matrix]\label{thm:elliptic-r}
For $\cA = \widehat{\fg}_k$ with $\fg$ simple and $k \neq -h^\vee$,
the KZ-normalized genus-$1$ two-point kernel equals the classical
Belavin elliptic kernel up to the standard level normalization:
\begin{equation}\label{eq:elliptic-rmatrix}
r_{\widehat{\fg}_k}^{(1)}(z,\tau)
\;=\;
\frac{1}{k+h^\vee}\,r^{\mathrm{ell}}_\fg(z,\tau).
\end{equation}
For $\fg = \fsl_2$, the Belavin $r$-matrix is
\begin{equation}\label{eq:sl2-elliptic}
r^{\mathrm{ell}}_{\fsl_2}(z,\tau)
\;=\;
\xi_\tau(z)\cdot\frac{H \otimes H}{2}
\;+\;
\phi_+(z,\tau)\,E \otimes F
\;+\;
\phi_-(z,\tau)\,F \otimes E,
\end{equation}
where $\{H,E,F\}$ is the standard basis of $\fsl_2$ and the
theta-function ratios are
\begin{equation}\label{eq:phi-pm}
\phi_\pm(z,\tau)
\;=\;
\frac{\theta_1'(0\mid\tau)\,
 \theta_1(z \pm \tfrac{1}{2}\mid\tau)}%
{\theta_1(z\mid\tau)\,\theta_1(\pm\tfrac{1}{2}\mid\tau)}.
\end{equation}
The full \(n\geq3\) KZB flatness statement uses the Felder dynamical
shift and the dynamical Yang--Baxter equation. It is not a consequence
of applying the ordinary rational CYBE to the displayed two-point
kernel.
\end{theorem}

\begin{proof}
For the Cartan channel: the Kac--Moody OPE with
$c_0 = \Omega/(k+h^\vee)$ gives, after genus-$1$ regularization,
$\xi_\tau(z)\cdot H \otimes H/2$ in the Cartan direction.

For the root channels: on $E_\tau$, the chiral fields $J^E$, $J^F$
associated to the positive and negative roots satisfy twisted boundary
conditions under $B$-cycle monodromy. The collision residue in the
root channels acquires the theta-function ratios
$\phi_\pm(z,\tau)$~\eqref{eq:phi-pm}, which are the unique
meromorphic functions on $\mathbb{C}$ with a simple pole at $z = 0$
matching $1/z$ and with $B$-cycle quasi-periodicity dictated by the
root lattice~\cite{Belavin81}.

The full \(n\)-point flatness statement belongs to the Felder dynamical
KZB lane. The two-point identification~\eqref{eq:elliptic-rmatrix}
follows from
Theorem~\ref{thm:elliptic-reg} combined with uniqueness of the
classical elliptic $r$-matrix~\cite{BelavinDrinfeld82}.
\end{proof}

\begin{remark}[Higher rank]\label{rem:higher-rank}
For $\fg = \fsl_N$ with $N \geq 3$, the Belavin elliptic $r$-matrix
is a matrix-valued meromorphic function on $E_\tau$ with poles at
$z = 0$ whose residue is the $\fsl_N$ Casimir. The collision residue
produces this $r$-matrix by the same mechanism: elliptic
regularization of the Cartan and root channels, with theta-function
ratios $\phi_\alpha(z,\tau)$ indexed by the roots $\alpha$. The
identification~\eqref{eq:elliptic-rmatrix} holds for all simple $\fg$.
\end{remark}


%%% ================================================================
%%% 5. FACE 7: THE ELLIPTIC GAUDIN MODEL
%%% ================================================================

\section{Face 7: the elliptic Gaudin model}\label{sec:gaudin}

At genus~$0$, Face~7 identifies the collision residue with the Gaudin
Hamiltonians
$H_i^{\mathrm{Gaudin}} = \sum_{j \neq i}\Omega_{ij}/(z_i - z_j)$.
At genus~$1$, the Gaudin model acquires elliptic structure.

\begin{theorem}[Face~7: elliptic Gaudin kernel]\label{thm:gaudin}
For $\cA = \widehat{\fg}_k$, the KZ-normalized two-point elliptic
Gaudin kernel on $\Conf_n(E_\tau)$ is
\begin{equation}\label{eq:elliptic-gaudin}
H_i^{\mathrm{ell}}
\;=\;
\frac{1}{k+h^\vee}\sum_{j \neq i}
\Omega_{ij}\,\xi_\tau(z_i - z_j)
\;=\;
\sum_{j \neq i}
\bigl[r_{\widehat{\fg}_k}^{(1)}(z_i - z_j,\tau)\bigr]_{ij}.
\end{equation}
After adding the Felder dynamical shift, this kernel is the
two-point sector of the flat KZB/Gaudin connection.
\end{theorem}

\begin{proof}
The identification of $H_i^{\mathrm{ell}}$ with the collision residue
is immediate from Theorem~\ref{thm:elliptic-reg}:
$r_{\widehat{\fg}_k}^{(1)}(z,\tau) = \Omega\,\xi_\tau(z)/(k+h^\vee)$
in KZ normalization. For \(n\geq3\), flatness of the elliptic Gaudin
connection requires the dynamical Cartan variable and Felder's
dynamical Yang--Baxter equation; the non-dynamical two-point kernel
displayed here is the affine Cartan/KZB comparison surface.
\end{proof}

\begin{remark}[Hitchin system]\label{rem:hitchin}
The elliptic Gaudin model is the genus-$1$ instance of the Hitchin
integrable system on $\Bun_G(E_\tau)$. The Hitchin Hamiltonians are
sections of a line bundle on the base of the Hitchin fibration; at
genus~$1$, the base is the affine space of elliptic spectral curves,
and the Hitchin Hamiltonians restrict to the elliptic Gaudin
Hamiltonians on the trivial $G$-bundle. At genus $g \geq 2$, the
collision residue on $\Sigma_g$ produces the Hitchin Hamiltonians on
$\Bun_G(\Sigma_g)$; this is the structural target of the holographic
construction~\cite{Lorgat26}.
\end{remark}


%%% ================================================================
%%% 6. CLASS M AT GENUS 1
%%% ================================================================

\section{Class $M$ at genus~$1$: higher-order elliptic
operators}\label{sec:class-m}

For affine Kac--Moody algebras (class~$L$), the collision expansion
terminates at $c_0$: the collision depth is $k_{\max} = 1$, and only
the \(\xi_\tau\)-term contributes. For class~$M$ algebras (Virasoro,
$\cW_N$), the OPE has poles of higher order, and the collision
expansion involves \(P_\tau\), \(P_\tau'\), and higher prime-form
derivatives.

\subsection{Virasoro at genus~$1$}

The Virasoro algebra has maximal OPE pole order $p_{\max} = 4$ (from
the $T(z)\,T(w)$ OPE). The bar propagator absorbs one order, giving
collision depth $k_{\max} = 3$. The genus-$0$ collision residue is
\begin{equation}\label{eq:vir-genus0}
r_{\Vir_c}(z) \;=\; \frac{c/2}{z^3} \;+\; \frac{2T}{z},
\end{equation}
with $c_0 = 2T$, $c_1 = 0$ (the $z^{-2}$ pole is absent from the
$r$-matrix by a parity argument: the $d\log$ extraction sends
$z^{-2n}$ to $z^{-(2n-1)}$), and $c_2 = c/2$.

\begin{theorem}[Virasoro genus-$1$ collision
residue]\label{thm:virasoro-g1}
The genus-$1$ collision residue of $\Vir_c$ is
\begin{equation}\label{eq:vir-genus1}
r_{\Vir_c}^{(1)}(z,\tau)
\;=\;
2T\cdot\xi_\tau(z) \;-\; \frac{c}{4}\,P_\tau'(z).
\end{equation}
This is an elliptic differential operator of order~$2$ on conformal
blocks on $E_\tau$: the \(\xi_\tau\)-term acts by multiplication,
while the \(P_\tau'\)-term acts as a second-order differential operator
through the stress-energy Ward identity.
\end{theorem}

\begin{proof}
The genus-$0$ collision expansion has $c_0 = 2T$, $c_1 = 0$,
$c_2 = c/2$. By Theorem~\ref{thm:elliptic-reg}:
\[
r_{\Vir_c}^{(1)}(z,\tau)
\;=\;
c_0\,\xi_\tau(z) + c_1\,P_\tau(z)
- \tfrac{1}{2}\,c_2\,P_\tau'(z).
\]
Since $c_1 = 0$, the \(P_\tau\)-term vanishes, and the result
is~\eqref{eq:vir-genus1}. The Laurent expansion near $z = 0$
confirms consistency: $2T/z - (c/4)\cdot(-2/z^3 + O(z))
= (c/2)/z^3 + 2T/z + O(1)$, matching $r_{\Vir_c}(z)$ at leading
order.

The \(P_\tau'\)-term acts as a second-order differential operator because
the Virasoro OPE at order $(z-w)^{-4}$ involves
$(c/2)\cdot\mathrm{id}$, which through the Ward identity produces a
second derivative $\partial_z^2$ on the correlator; the elliptic
regularization replaces $1/z^4$ by $-P_\tau'(z)/2$.
\end{proof}

\begin{remark}[The $c_1 = 0$ vanishing]\label{rem:c1-vanishing}
The vanishing of $c_1$ is a parity phenomenon: the OPE mode
$T_{(2)}T = \partial T$ is a total derivative, and its contribution to
the $r$-matrix vanishes after $d\log$-extraction. In the elliptic
regularization, this means the \(P_\tau(z)\)-term is absent. The
first non-trivial correction beyond \(\xi_\tau\) involves \(P_\tau'\),
not \(P_\tau\).
\end{remark}


\subsection{$\cW_N$ at genus~$1$}

The $\cW_N$ algebra has generators $W_2 = T, W_3, \ldots, W_N$ of
conformal weights $2,3,\ldots,N$. The maximal OPE pole order among
generating fields is $p_{\max} = 2N$ (from the $W_N$-$W_N$ OPE),
giving collision depth $k_{\max} = 2N-1$.

\begin{theorem}[$\cW_N$ genus-$1$ collision
residue]\label{thm:wn-g1}
The genus-$1$ collision residue of $\cW_N$ involves the prime-form
kernel and the derivatives of \(P_\tau\) up to \(P_\tau^{(2N-3)}\):
\begin{equation}\label{eq:wn-genus1}
r_{\cW_N}^{(1)}(z,\tau)
\;=\;
c_0\,\xi_\tau(z)
\;+\;
\sum_{\substack{n=1 \\ c_n \neq 0}}^{2N-2}
\frac{(-1)^{n-1}}{n!}\,c_n\cdot P_\tau^{(n-1)}(z),
\end{equation}
where $c_n$ are the OPE-mode collision coefficients. The
highest-order term is an elliptic differential operator of order
$2N-2$. For $\cW_3$: the collision expansion involves
\(\xi_\tau,P_\tau,P_\tau',P_\tau'',P_\tau'''\), producing fourth-order elliptic
differential operators.
\end{theorem}

\begin{proof}
Direct specialization of Theorem~\ref{thm:elliptic-reg} to the
$\cW_N$ OPE data with $p_{\max} = 2N$: the collision depth
$k_{\max} = 2N-1$ means the genus-$0$ expansion
$r_{\cW_N}(z) = \sum_{n=0}^{2N-2}c_n/z^{n+1}$ has $2N-1$ terms,
each regularized by \(1/z\mapsto\xi_\tau(z)\) and
\(1/z^{n+1}\mapsto (-1)^{n-1}P_\tau^{(n-1)}(z)/n!\) for \(n \geq 1\).
\end{proof}

\begin{remark}[Outside the affine KZB lane]\label{rem:genuinely-new}
The elliptic differential operators produced by the class~$M$ collision
residue lie outside the affine Hitchin--KZB \(r\)-matrix lane. The
Hitchin system on $\Bun_G(E_\tau)$ sees the affine Casimir channel.
The KZB connection sees the prime-form kernel and its first polar
derivative in the standard affine sector. The class~$M$ collision
residue at depth $\geq 3$ produces operators involving
\(P_\tau',P_\tau'',\ldots,P_\tau^{(2N-3)}\) from higher OPE poles
specific to $\cW$-algebras.
\end{remark}


%%% ================================================================
%%% 7. DEGENERATION TO GENUS 0
%%% ================================================================

\section{Degeneration to genus~$0$}\label{sec:degeneration}

\begin{theorem}[Degeneration]\label{thm:degeneration}
As $\mathrm{Im}(\tau) \to \infty$ (equivalently,
$q = e^{2\pi i\tau} \to 0$), all seven genus-$1$ faces degenerate to
their genus-$0$ counterparts:
\begin{enumerate}[label=\textup{(\roman*)}]
\item \textup{(Propagator.)}
\(\xi_\tau(z) \to \pi\cot(\pi z) \to 1/z + O(z)\).

\item \textup{(Prime-form derivatives.)}
\(P_\tau(z) \to \pi^2/\sin^2(\pi z) \to 1/z^2 + O(1)\), and
\(P_\tau^{(m)}(z) \to (-1)^m(m+1)!/z^{m+2} + O(z^{-m})\).

\item \textup{(KZB $\to$ KZ.)} The $dz$-component degenerates to
the KZ connection: \(\xi_\tau(z_{ij}) \to 1/z_{ij}\). The
$d\tau$-component has no genus-$0$ counterpart.

\item \textup{(Elliptic $\to$ rational $r$-matrix.)}
In trace form,
\[
r^{\mathrm{KM}}_{\mathrm{tr}}(z)=k\,\Omega_{\mathrm{tr}}/z.
\]
The KZ-normalized connection uses
\[
r^{\mathrm{KZ}}(z)=\Omega_{\mathrm{KZ}}/((k+h^\vee)z),
\qquad
\Omega_{\mathrm{KZ}}=2h^\vee\Omega_{\mathrm{tr}}.
\]
The elliptic kernel degenerates through the trigonometric kernel to
the corresponding rational kernel in the chosen normalization. These
are conventionally related normalizations, not an identity as rational
functions of~$k$.

\item \textup{(Elliptic $\to$ rational Gaudin.)}
$H_i^{\mathrm{ell}} \to H_i^{\mathrm{Gaudin}} = \sum_{j \neq i}
\Omega_{ij}/(z_i - z_j)$.
\end{enumerate}
\end{theorem}

\begin{proof}
As \(q\to0\), the prime-form kernels have the standard degenerations.
The logarithmic derivative satisfies
\(\xi_\tau(z)=\pi\cot(\pi z)+O(q)\), giving \(1/z+O(z)\) for
\(|z|<1\). Its derivative gives
\(P_\tau(z)=\pi^2/\sin^2(\pi z)+O(q)\), hence \(1/z^2+O(1)\). Higher
derivatives degenerate analogously.

The KZB degeneration: with \(\xi_\tau(z) \to 1/z\), the
$dz$-component~\eqref{eq:kzb-z} becomes
$\sum_{j \neq i}\Omega_{ij}/((k+h^\vee)z_{ij})\,dz_{ij}$, which is
the KZ connection.

The elliptic $r$-matrix degeneration: as $q \to 0$,
$\phi_\pm(z,\tau) \to \pi/\sin(\pi z) + O(q)$, recovering the
trigonometric $r$-matrix. A further limit
$z \mapsto z/L$, $L \to \infty$ gives $\pi/\sin(\pi z/L) \to L/z$,
recovering the rational $r$-matrix
$\Omega/((k+h^\vee)z)$ in the KZ normalization.
Both limits are continuous in the topology of formal Laurent series.
\end{proof}

The degeneration theorem provides the consistency check between genus-$0$
and genus-$1$: every genus-$1$ face is a continuous deformation of the
corresponding genus-$0$ face, with deformation parameter $q$. At
$q = 0$, the elliptic curve degenerates to the nodal rational curve,
and all seven identifications collapse to their genus-$0$ ancestors.


%%% ================================================================
%%% 8. COMPUTATIONAL VERIFICATION
%%% ================================================================

\section{Computational verification}\label{sec:compute}

The genus-$1$ seven-face checks are implemented by the compute
engine \texttt{theorem\_genus1\_seven\_faces\_engine.py}, which
implements the following tests.

\begin{enumerate}[label=(\arabic*)]
\item \textbf{Elliptic regularization.} For each standard family
$\cA$ in the landscape census, verify that the genus-$1$ collision
residue~\eqref{eq:elliptic-expansion} matches the elliptic
regularization of the genus-$0$ collision expansion at \(O(q^0)\)
within the stated numerical tolerance.

\item \textbf{KZB identification.} For
$\cA = \widehat{\fsl}_2$ at levels $k = 1,2,3,5,10$, verify that the
$dz$-component of the KZB connection equals the collision residue, and
that the $d\tau$-component equals $-\partial_z r_\cA^{(1)}$.

\item \textbf{Dynamical closure.} Check the two-point and Cartan
sectors of the affine elliptic kernel. The full \(n\geq3\) flatness
statement is the Felder dynamical Yang--Baxter closure and is not
claimed by this ordinary-kernel test.

\item \textbf{Gaudin kernel.} Check that the two-point and Cartan
elliptic Gaudin kernels agree with the KZB kernel in the chosen
normalization. Full \(n\geq3\) commutativity belongs to the dynamical
KZB/Gaudin lane.

\item \textbf{Degeneration.} Verify that
$\|r_\cA^{(1)}(z,\tau) - r_\cA(z)\|/\|r_\cA(z)\| < 10^{-6}$
for $\mathrm{Im}(\tau) \geq 12$ across all standard families.

\item \textbf{Class~$M$ structure.} For $\Vir_c$ at
$c = 1, 1/2, 25, 26$, verify~\eqref{eq:vir-genus1} and confirm the
$c_1 = 0$ vanishing. For $\cW_3$ at $c = 2, 50, 98$,
verify~\eqref{eq:wn-genus1} and confirm terms through \(P_\tau'''\).
\end{enumerate}


%%% ================================================================
%%% 9. VERLINDE FORMULA
%%% ================================================================

\section{The Verlinde formula from ordered chiral
homology}\label{sec:verlinde}

At integer level $k \in \mathbb{Z}_{\geq 1}$, the ordered chiral
homology of $\widehat{\fg}_k$ on $E_\tau$ recovers the Verlinde
partition function. The connection to the genus-$1$ seven-face
identification is
direct: the KZB connection of Theorem~\ref{thm:kzb} has finite
monodromy at integer level, and the trace of the $B$-cycle monodromy
over the space of conformal blocks recovers the Verlinde formula.

\begin{proposition}[Verlinde from genus-$1$ collision
residue]\label{prop:verlinde-g1}
For $\cA = \widehat{\fg}_k$ at integer level $k \geq 1$, the $B$-cycle
shift of the KZ-normalized prime-form kernel is
\begin{equation}\label{eq:b-monodromy}
\mathcal{M}_{B}^{\mathrm{ker}}
\;=\;
\exp\Bigl(-\frac{2\pi i}{k+h^\vee}\,\Omega\Bigr)
\;=\;
\exp(-2\pi i\,c_0),
\end{equation}
where $c_0 = \Omega/(k+h^\vee)$ is the degree-$0$ collision
coefficient. The quantum group parameter is
\(q_{\mathrm{DK}}=\exp(\pi i/(k+h^\vee))\), with full-cycle factor
\(q_{\mathrm{DK}}^2=\exp(2\pi i/(k+h^\vee))\). The Verlinde partition function
at genus~$1$ is
\begin{equation}\label{eq:verlinde-g1}
Z_1(\widehat{\fg}_k)
\;=\;
\sum_{\lambda \in P_+^k}
|S_{0\lambda}|^{0}
\;=\;
|P_+^k|,
\end{equation}
where $P_+^k$ is the set of integrable highest weights at level $k$
and $S_{0\lambda}$ are the modular $S$-matrix entries. At higher genus
$g$, the Verlinde formula
$Z_g = \sum_\lambda S_{0\lambda}^{2-2g}$ is recovered from the
ordered chiral homology via handle attachment and separating
factorization of the ordered bar complex.
\end{proposition}

\begin{remark}[The $B$-cycle monodromy and the quantum
group]\label{rem:quantum-group}
The quasi-periodicity~\eqref{eq:B-cycle-monodromy} of the genus-$1$
collision residue is the geometric origin of the quantum group
parameter. For affine $\widehat{\fsl}_{2,k}$: the $B$-cycle monodromy
eigenvalues are written in the Drinfeld--Kohno normalization with
\(q_{\mathrm{DK}}=e^{\pi i/(k+2)}\), and the full cycle contributes
\(q_{\mathrm{DK}}^2\). The passage from the
$E_1$ collision residue $r^{(1)}(z,\tau)$ to the quantum group
\(U_{q_{\mathrm{DK}}}(\fg)\) is the $B$-cycle monodromy of the ordered bar complex:
the spectral parameter $z$ (genus-$0$ data) becomes the quantum
group parameter through exponentiation of the
quasi-period.
\end{remark}


%%% ================================================================
%%% 10. THE CRITICAL LEVEL JUMP
%%% ================================================================

\section{The critical level}\label{sec:critical-level}

The genus-$1$ seven-face identification has a distinguished boundary:
the critical level $k = -h^\vee$, where the Sugawara construction
degenerates. The collision residue and the KZB connection undergo a
qualitative change at this value.

\begin{theorem}[Critical level jump at
genus~$1$]\label{thm:critical-level-g1}
For $\cA = \widehat{\fg}_k$ at the critical level $k = -h^\vee$:
\begin{enumerate}[label=\textup{(\roman*)}]
\item The shadow invariant vanishes:
$\kappa(\widehat{\fg}_{-h^\vee}) = 0$. The KZB-normalized kernel is
singular at this level, so the non-critical monodromy formula has no
regular value there.
\item The degree-$2$ ordered chiral homology on $V \otimes V$ doubles:
$\dim H^1_{\mathrm{ord}, 2}$ jumps from $4$ (generic level) to $8$
(critical level) for $\fg = \fsl_2$ \textup{(}finite computation,
$\chi_{\mathrm{ord}, 2} = -4$ conserved\textup{)}. This is distinct from
the full symmetric bar cohomology at critical level, addressed in (iii).
\item Under the completed critical bar--Ext comparison, the critical
bar--Ext groups identify with differential forms on the space of
\(\fg^\vee\)-opers:
\[
H^n_{\mathrm{bExt}}(\widehat{\fg}_{-h^\vee})
\cong \Omega^n(\mathrm{Op}_{\fg^\vee}(D)).
\]
This is a completed critical statement, not an equality for the raw
symmetric bar cohomology.
\item The genus-$1$ collision residue degenerates: the KZ normalization
$r^{(1)}(z,\tau) = \Omega\,\xi_\tau(z)/(k+h^\vee)$ has a pole at
$k = -h^\vee$, and the KZB connection becomes singular. The critical
level requires a separate regularization via the Feigin--Frenkel
center~\cite{FeiginFrenkel94}.
\end{enumerate}
\end{theorem}

\begin{remark}[Koszulness at critical level]\label{rem:koszul-critical}
At generic level $k \neq -h^\vee$, the affine Kac--Moody algebra
$\widehat{\fg}_k$ is chirally Koszul, and the seven-face
identification applies without modification. At the critical level
$k = -h^\vee$,
the KZ/KZB normalization and Sugawara topologisation fail, and the
Feigin--Frenkel centre appears. This does not by itself negate PBW
Koszulness of the universal affine algebra; it moves the statement to
the separate critical-level bar-cohomology lane. The seven KZ faces
degenerate: the Gaudin Hamiltonians become the Feigin--Frenkel higher
Gaudin
Hamiltonians, and the $r$-matrix limit requires the Langlands dual
group ${}^L G$ via the geometric Langlands correspondence. The
critical level is a boundary of the non-critical seven-face
identification.
\end{remark}

\begin{remark}[Spectral sequence at critical
level]\label{rem:spectral-critical}
The passage from generic to critical level is controlled by a spectral
sequence with $d_1 = \lambda \cdot [\delta]$, where
$\lambda = k + h^\vee$ is the level shift and $\delta$ is the boundary
class. At $\lambda = 0$ (critical level), $d_1$ vanishes, and the
spectral sequence degenerates at $E_1$: the bar cohomology jumps
because the differential that would kill the extra classes is absent.
This is the mechanism by which the degree-$2$ ordered chiral homology
$H^1_{\mathrm{ord}, 2}$ doubles from $4$ to $8$ for $\fsl_2$. For the
completed critical bar--Ext cohomology at critical level, the
\(d_1\)-vanishing
lets the entire de~Rham algebra $\Omega^*(\mathrm{Op}_{\fsl_2}(D))$
appear in completed critical bar--Ext cohomology, which is
infinite-dimensional.
\end{remark}


%%% ================================================================
%%% REFERENCES
%%% ================================================================

\begin{thebibliography}{10}

\bibitem{Belavin81}
A.~A.~Belavin,
\emph{Dynamical symmetry of integrable quantum systems},
Nucl.\ Phys.\ B~\textbf{180} (1981), 189--200.

\bibitem{BelavinDrinfeld82}
A.~A.~Belavin and V.~G.~Drinfeld,
\emph{Solutions of the classical Yang--Baxter equation for simple Lie
algebras},
Funct.\ Anal.\ Appl.\ \textbf{16} (1982), 159--180.

\bibitem{Bernard88}
D.~Bernard,
\emph{On the Wess--Zumino--Witten models on the torus},
Nucl.\ Phys.\ B~\textbf{303} (1988), 77--93.

\bibitem{DNP25}
T.~Dimofte, W.~Niu, and A.~Py,
\emph{Line operators in $3$d holomorphic QFT: meromorphic tensor
categories and dg-shifted Yangians},
arXiv:2508.11749, 2025.

\bibitem{FeiginFrenkel94}
B.~Feigin, E.~Frenkel, and N.~Reshetikhin,
\emph{Gaudin model, Bethe ansatz and critical level},
Comm.\ Math.\ Phys.\ \textbf{166} (1994), 27--62.

\bibitem{KhanZeng25}
A.~Z.~Khan and K.~Zeng,
\emph{Poisson vertex algebras and three-dimensional gauge theory},
arXiv:2502.13227, 2025.

\bibitem{KZB95}
D.~Bernard,
\emph{On the Wess--Zumino--Witten models on Riemann surfaces},
Comm.\ Math.\ Phys.\ \textbf{124} (1989), 315--340;
V.~G.~Knizhnik and A.~B.~Zamolodchikov,
\emph{Current algebra and Wess--Zumino model in two dimensions},
Nucl.\ Phys.\ B~\textbf{247} (1984), 83--103.

\bibitem{Lorgat26}
R.~Lorgat,
\emph{Modular Koszul Duality},
preprint, 2026.

\bibitem{STS83}
M.~A.~Semenov-Tian-Shansky,
\emph{What is a classical $r$-matrix?},
Funct.\ Anal.\ Appl.\ \textbf{17} (1983), 259--272.

\end{thebibliography}

\end{document}
